\documentclass[a4paper, 12pt]{article}

\input{preamble}

\begin{document}
% This sets the default language for the listings package.
\lstset{language=Python}
\title{
\sffamily
\vspace{-1.7cm}
{\huge\textbf{Project 2b -- Bayesian Optimization of Ad-Atom Energy Landscapes}} \\[-0.4cm]
}
\author{
    \vspace{-0.4cm}
    \begin{minipage}[t]{0.45\textwidth}
\centering
Linus Brink\\
\href{mailto:brinkl@chalmers.se}{brinkl@chalmers.se}
\end{minipage}%
\hfill
\begin{minipage}[t]{0.45\textwidth}
\centering
Oscar Stommendal\\
\href{mailto:oscarsto@chalmers.se}{oscarsto@chalmers.se}
\end{minipage}
}
\date{\today}
\maketitle
% \begin{center}
%     \textbf{Abstract}
% \end{center}
\vspace{-1.5cm}
\section{Introduction}
Despite our never-ending love for physics, we can not deny the mental fatigue that comes with simulating and understanding complex systems. As the systems grow in complexity, so does the computational cost of simulating them. This is especially true in the field of material science, where we often want to understand the properties of materials at the atomic scale. To overcome this challenge, we can use surrogate models that approximate the behavior of the system based on a limited number of simulations. One such approach is Bayesian Optimization (BO), which combines probabilistic modeling with optimization techniques to efficiently explore the parameter space. In this project, we apply BO to model the energy landscape of an ad-atom on a gold-surface. By using Gaussian Processes (GPs) as our surrogate model, we aim to predict the energy of the ad-atom at unobserved locations on the surface, while minimizing the number of expensive energy calculations required. This report details our methodology, results, and insights gained from applying BO to this problem.

\section{Theory}
For an ad-atom on a surface, its energy $E$ can be expressed as 
\begin{equation}
E = E_{\text{ad}} - E_{\text{surface}},
\end{equation} 
where $E_\text{surface}$ is the energy of the bare surface and $E_\text{ad}$ is the energy of the combined system. If we consider the surface to be completely rigid, the energy landscape of the ad-atom can be described as a function of its position $(x, y, z)$ on the surface. By fixing the height $z$ to its local minimum given $(x, y)$, we can reduce the problem to a two-dimensional energy landscape $E(x, y)$. By combining GPs with BO, we can efficiently model this energy landscape and predict the energy at unobserved locations, while minimizing the number of expensive energy calculations required.

\section{Methodology}
In this project, we consider a gold (Au) structure with Miller indices (433) and an ad-atom of gold placed on its surface. The energy of the ad-atom is calculated using the Atomic Simulation Environment (ASE) package using the Effective Medium Theory (EMT) potential \cite{ase-paper}. As outlined in the introduction, we model the energy landscape of the ad-atom as a function of its position $(x, y)$ on the surface, while keeping its height $z$ fixed at the local minimum for each $(x, y)$ pair, i.e., 
\begin{equation}
E(x, y) = \min_z E(x, y, z).
\end{equation}
To perform the minimization over $z$, we use a scalar optimization method from the \texttt{SciPy} library to find the local minimum of the energy with respect to $z$ for each fixed $(x, y)$ pair \cite{2020SciPy-NMeth}. Furthermore, when searching for the optimal ad-atom position, the energy landscape will repeat itself due to the periodicity of the crystal structure. Therefore, we only need to consider a single unit cell of the surface, i.e., the domain
\begin{equation}
    0 < x < 16.65653 \quad \land \quad 0 < y < 2.884996.
\end{equation}
We used a step size of 0.1 \AA{} in both $x$ and $y$ directions when creating the grid of possible ad-atom positions, resulting in a (29 x 167) grid. 

For an initial picture of the energy landscape, it was plotted over the entire grid. From this, the approximate position of the global minimum could be identified for further comparison and analysis. Next, we applied local search using \texttt{SciPy}, specifically L-BFGS-B, to investigate the performance of regular gradient-based optimization on this energy landscape. The algorithm was initialized at 500 random positions within the unit cell, and the number of times it converged to the global minimum was recorded. 

\section{Results}
Figure~X shows the potential energy surface $E(x,y)$ of the Au ad-atom evaluated on the $29 \times 167$ grid over the primitive cell. 
The energy varies between roughly $7.05$ and $8.1$~eV, with several local minimas and maximas in a reoccuring pattern. The pattern breaks for $x \lesssim 5$, where the global extremes appear. The global minima is located at $(3.31, 0.72)$ with an energy of $E_\text{min} = 7.0783$\,eV. The region around the global minimum is rather small as it is narrow in $x$ and moderately extended in $y$. 

\textcolor{red}{Insert Figure 1 of PES for task 1}

Something about task 2...

Figure~Y shows, for each value of $\beta$, the energy of the sampled position $(x,y)$ at each Bayesian optimisation iteration. The dashed horizontal line marks the global minimum energy obtained from the full PES. For all $\beta$ values that manage to locate the global minimum, the model reaches this energy level relatively early in the run, but then leaves the minimum again and explores other regions of the energy landscape before eventually returning to the global minimum and converging there.

Something about which one is "best" and how beta affects the convergence...

\textcolor{red}{Insert Figure of sampling for different betas from task 3}

Figure~Z illustrates the final BO GP model for the optimal $\beta$ that converged to the global minimum the fastest. The left panel shows the predicted mean energy over the domain together with the sampled positions and the known global minimum. Most sampled positions lie in low energy regions, with a dense cluster in the well that contains the global minimum. The middle panel shows the predictive standard deviation (uncertainty) $\sigma$. The uncertainty is small in areas that contain many sampled positions, generally around the minimas (especially the global minimum), and larger in regions that have not been visited. The right panel displays the value of the LCB acquisition function for the same model. It basically shows the opposite of the ucnertainty. High acquisition values occur in the low energy region surrounding the minimas where the data has been, and low where it has not.

\textcolor{red}{Insert Figure over optimal beta + variance + aquisition function + data points from task 3}

Figure~Å shows the root mean square error (RMSE) of the general GP model evaluated on the full PES as a function of the number of training samples. The RMSE starts at about $0.30$\,eV for very small training sets and decreases as more samples are added, where the most rapid improvement occurs during the first $\sim 80$ samples. After $\sim 120$ the curve has basically flattend around $0.02$–$0.03$\,eV until 200 training samples, which indicates that the model has captured the main structure of the PES.

\textcolor{red}{Insert Figure with the RMSE from task 4}

Figure~Ä compares the energy along the linear transition path between the global minimum ($\lambda = 0$) and the local energy minimum around $(11, 2.1)$ ($\lambda = 1$) between the BO and general GP with the target energy along the path computed with EMT. The BO GP model starts identically to the target with basically no uncertainty around the global minimum, but shows a higher energy for both the first local maximum and minimum, followed by combining the next two maxima and the minimum into one larger maximum with a large uncertainty. The BO GP model is very good close to the global minimum, but smoothes the rest of the path. The general GP is slightly worse at the global minimum but keeps a very low differential to the target with a lower uncertainty throughout the path, being a lot better overall.

\textcolor{red}{Insert Figure with the two models and the actual energy from task 4}



\section{Discussion}
The appearance of the PES in Figure~X explains why the optimization problem can be demanding. Many narrow minima are spread over the unit cell and the global minimum does not stand out clearly by shape alone. Only a small fraction of the domain lies close to the global minimum, which means that a method that only follows local gradients is expected to reach it from relatively few starting positions.

The Bayesian optimization traces in Figure~Y show that the GP based search can still identify the global minimum efficiently. Write about the betas that worked...

For the values of $\beta$ that succeed, the sampled energy reaches the global minimum early, then moves away again while other parts of the PES are explored, and then finally returns to converge at the global minimum. This behavior matches the idea behind the LCB acquisition function, where low predicted energy and high uncertainty are both rewarded. After the first visit the minimum looks very good in terms of mean energy, but unexplored wells still have enough uncertainty to be tested before the search settles. The dependence on $\beta$ reflects the balance between these two effects. A smaller $\beta$ favors exploitation, increasing the chance to stay around already discovered low-energy regions and thus shortening the exploration time, but it also increases the risk of remaining trapped in a local well. A larger $\beta$ instead favors exploration and improves the chance of visiting the global minimum at least once, but can give too much weight to uncertainty and spend many evaluations in high-energy regions that are unlikely to contain the global minimum. Hopefully this shows somewhat in the figures...

The spatial plots in Figure~Z help to interpret this behavior. The predictive uncertainty is low near sampled positions and especially around minima that have been revisited, while it remains high in regions that were never chosen by the acquisition function. High acquisition values appear where low predicted energy and nonzero uncertainty coincide, which is mainly in and around the wells that contain sampled points. Regions that are clearly high in energy receive low acquisition values even if they remain unsampled. This pattern shows that the combined GP and LCB scheme focuses evaluations on low-energy regions that already contain samples, while largely ignoring unsampled regions that are clearly high in energy.

The general GP model in Task~4 addresses a different goal, namely to approximate the full PES rather than only to locate the global minimum. The RMSE in Figure~Å decreases rapidly for the first 80 samples and then levels off around a few hundredths of an eV after about 120 samples, which suggests that the main structure of the PES has been captured and that additional data mainly refine the model. This means that the optimal amount of training samples likely is around 120, maybe even as early as 80. 

The comparison of the transition path in Figure~Ä supports this view. The BO based GP is very accurate close to the global minimum, where most of its samples lie, while the general GP, trained on a broader set of points, follows the reference energy more closely along the entire path.



\printbibliography

\end{document}